% 来源：OEBPS/Text/01_39.xhtml（OPF spine 顺序；仅第一册）
\cleardoublepage
\BookEntry[弗兰克·帕特诺伊《诚信的背后》]{弗兰克·帕特诺伊《诚信的背后》}

\BookCover{images/书单24诚信的背后.jpg}{[美] 弗兰克·帕特诺伊著| 邵琰译| 当代中国出版社2008年版}

\begin{BookSummaryBox}
{\centering \textbf{揭露华尔街衍生品黑幕\\ 第一本关于衍生品的国际畅销书}\par}
\end{BookSummaryBox}

\subsubsection*{这本书谈什么？}

《诚信的背后》是弗兰克·帕特诺伊以自身经历写成的一部职场回忆录，揭露了以摩根士丹利为代表的投资银行，如何在金融衍生品贸易上进行所谓的“大胆创新”，就算损害客户利益也在所不惜，已沦为巨型赚钱机器的事实。这也是第一本关于金融衍生品的国际畅销书。
\par

1993年，年轻的弗兰克·帕特诺伊从耶鲁大学法学院毕业，进入了第一波士顿，次年转投著名投资银行摩根士丹利，成了世界上最能赚钱的一群人的一员。1995年，他离开摩根士丹利，以实实在在的亲身经历写成《诚信的背后》，揭开了华尔街金融创新与圈钱游戏的真相。在华尔街，如果有人说你是个好人，他的言外之意就是你是个笨蛋。对于大多数金融产品（如期货、外汇、衍生物等）而言，买卖双方的盈亏之和等于零。也就是说一方的利润就是另一方的损失。书中提到的许多美国著名公司都是这些利润背后的牺牲品。
\par

\subsubsection*{芒格为什么谈起这本书？}

芒格谈到对金融机构和衍生品的警告时，最后考虑到衍生品破坏的可能性，他提到了这本书：
\par

\begin{quotation}
在美国，绝大多数大型公司无法甩掉衍生品账本。对于伯克希尔这种已经很有钱的公司而言，去从事衍生品业务是很神经的想法。对大银行来说，衍生品是一桩疯狂的生意。
\par

如果你是个头脑正常的人，曾经花一个月时间认真钻研大型的衍生品业务，你肯定会感到恶心。你会觉得它是刘易斯·卡洛尔（《爱丽丝漫游奇境记》的作者）。你会觉得它是“疯帽匠的茶话会”。这些人虚假的精确更是让人难以置信。他们把最糟糕的经济学教授奉为神明。去看《诚信的背后：华尔街圈钱游戏的真相》（又名《泥鸽靶》）那本书，作者是做过衍生品交易员的法学教授弗兰克·帕特诺伊，他揭露了华尔街规模最大、声望最佳的公司之一的衍生品交易黑幕。
\par
\end{quotation}

\subsubsection*{《诚信的背后》作者弗兰克·帕特诺伊如是说}

“我在摩根士丹利衍生产品部的面试进行得很顺利，尽管过程中我尽量不让自己显得过于急切。部门的名称是DPG，‘衍生产品部’（Derivative Products Group）的缩写。面试时一个经理咕哝了几句‘一流的业务，一流的文化’之类的话，另一个则说他确信衍生产品部的发展将很大程度上来自新兴市场，也就是我所在的领域。他许诺说如果我加入摩根士丹利，他会邀请我参加衍生产品部门的重要聚会——下一次活动的时间在1994年4月，名字叫作‘F.I.A.S.C.O.’——泥鸽靶射击比赛。我听说过这个名字，而且很想参加，但我知道当一个投资银行的经纪人向你许诺什么事时，这件事实际发生的可能性几乎不存在。”
\par
